\chapter{Static Timing Analysis Engine}
\label{ch:sta}

This chapter documents the implementation of the theory reviewed in
\cref{sec:bg-sta}, as it appears in \code{analysis/sta.py}. The module's
only two public entry points are

\begin{lstlisting}[style=pythonstyle, numbers=none]
def analyze_timing(circuit: Circuit,
                   gate_delays: Mapping[str, RiseFall] | None = None
                   ) -> TimingAnalysisResult: ...

def analyze_timing_metrics(circuit: Circuit,
                   gate_delays: Mapping[str, RiseFall] | None = None
                   ) -> TimingMetrics: ...
\end{lstlisting}

\code{analyze\_timing} is the full analysis used everywhere in this report
(\cref{ch:results}) and returns a \code{TimingAnalysisResult} --- a frozen
dataclass carrying \code{critical\_paths}, \code{wns}, \code{tns},
\code{circuit\_delay}, and the full \code{arrival\_times} /
\code{required\_times} / \code{transition\_slacks} tables (each a mapping
from net name to a \code{RiseFall} pair), plus a
\code{node\_slack(net\_name)} convenience method.
\code{analyze\_timing\_metrics} is a deliberately cheaper sibling, used by
the optimizer's inner evaluation loop (\cref{sec:opt-eval}) whenever the
full critical-path list is not needed for a particular candidate check. Its
speedup is not merely skipping the reporting of unneeded data: since
$\mathrm{WNS}$/$\mathrm{TNS}$ (\cref{eq:wns-tns}) only ever need each
primary output's \emph{own} configured required time
(\code{circuit.config.output\_required(net\_name)}) compared against that
output's forward-computed arrival time, \code{analyze\_timing\_metrics}
never performs the full backward propagation of \cref{eq:backward-sta}
through the circuit's internal nodes at all --- it calls only the lean
arrival-time pipeline (\code{\_compute\_arrival\_values} /
\code{\_maximum\_gate\_arrival}, \cref{sec:sta-forward}), which does not
even retain the tie-tracking predecessor information \code{analyze\_timing}
needs for path reconstruction, and computes each output's slack directly
against its own configured requirement.

Both public functions accept an optional \code{gate\_delays} override,
which the Monte Carlo module (\cref{ch:monte-carlo}) uses to inject a
sample's perturbed delay set without needing a variant analysis path of its
own. Internally, both are thin wrappers around a shared private pipeline in
\code{analysis/sta.py}, whose exact signatures are reproduced below since
they pin down what state flows between the forward pass, the
backward pass, and path extraction:

\begin{lstlisting}[style=pythonstyle, numbers=none]
def _run_sta(circuit: Circuit, gate_delays: Mapping[str, RiseFall]
             ) -> tuple[dict[str, RiseFall],   # arrival_times
                        dict[str, RiseFall],   # required_times
                        PredecessorTable,      # controlling arcs
                        dict[str, RiseFall]]:  # transition_slacks

def _compute_arrival_times(circuit, gate_delays
    ) -> tuple[dict[str, RiseFall], PredecessorTable]: ...   # Sec. 5.3
def _compute_required_times(circuit, gate_delays) -> dict[str, RiseFall]: ...     # Sec. 5.4
def _compute_transition_slacks(arrival_times, required_times
        ) -> dict[str, RiseFall]: ...                    # Sec. 5.5
def _find_critical_paths(circuit, predecessors, transition_slacks
        ) -> list[CriticalPath]: ...                     # Sec. 5.6
def _trace_paths_backward(circuit, current_net, current_transition,
        output_net, reversed_steps, reversed_transitions, endpoint_slack, found_paths, predecessors) -> None: ...    # Sec. 5.6
\end{lstlisting}

with \code{\_compute\_circuit\_delay} and \code{\_compute\_output\_slack\_metrics}
deriving \code{circuit \_delay}, \code{wns}, and \code{tns} from the
completed arrival/slack tables, and \code{\_timing\_result} assembling the
final \code{TimingAnalysisResult}. The type
\code{PredecessorTable = Mapping[tuple[str, Transition],
tuple[PredecessorArc, ...]]} is itself worth noting: it is keyed by
\emph{every} (net, transition) pair to a \emph{tuple} of controlling
predecessor arcs, not a single one --- this is precisely the data structure
that makes exhaustive tie recording (\cref{sec:sta-forward}) possible,
since \code{\_compute\_arrival\_times} records every arc attaining the
maximum in \cref{eq:forward-sta}, and \code{\_trace\_paths\_backward}
later branches over exactly that stored tuple rather than recomputing or
re-deciding predecessor selection from scratch.

The engine executes in five ordered steps for a given, fixed gate-sizing
assignment: (1) load and delay computation (\cref{sec:sta-load}), performed
eagerly by \code{Circuit.\_\_init\_\_} itself rather than by
\code{analyze\_timing} (\cref{sec:arch-principles}); (2) forward propagation
(\cref{sec:sta-forward}); (3) backward propagation
(\cref{sec:sta-backward}); (4) slack, WNS, and TNS computation
(\cref{sec:sta-slack}); and (5) critical-path extraction
(\cref{sec:sta-paths}). Steps (2) and (3) both consume the results of step
(1) but not of each other, and are logically independent of one another
aside from that shared dependency.

\section{Circuit Graph Construction and Topological Order}
\label{sec:sta-graph}

Once a netlist has passed validation (\cref{ch:validation}), it is
represented as an object graph rather than a collection of string-indexed
lookup tables: each net object holds a direct reference to the gate object
that drives it and a list of the \code{(pin\_index, gate)} pairs it loads,
and each gate object holds direct references to its input net objects and
its output net object. This means that finding ``who drives net $X$'' or
``what does gate $Y$ fan out to'' is an $O(1)$ pointer dereference rather
than a linear scan over every gate in the circuit, which matters
materially at the scale of the optimizer's inner loop
(\cref{ch:optimization}), where the same queries are issued on the order of
hundreds of times per circuit.

The topological order is computed once, at \code{Circuit} construction, by
the module-level helper \code{\_make\_topological\_order(netlist, gates) ->
tuple[tuple[Gate, ...], ...]} (\code{domain/circuit.py}) using Kahn's
algorithm, and is stored not as a flat list but as a sequence of
\emph{levels} --- exactly the nested-tuple return type above: gates are
grouped by their longest-path depth from any
primary input, so that all gates in level $\ell$ depend only on gates in
levels $0, \ldots, \ell-1$, and this is the \code{topological\_order}
field consumed directly by \code{\_compute\_arrival\_times}'s outer
\code{for level in circuit.topological\_order:\allowbreak\ for gate in
level:} loop (\cref{sec:sta-forward}). Flattening the levels in order recovers an
ordinary topological order; the level structure itself is not required for
correctness of the forward/backward passes below (which only need \emph{a}
topological order, not a minimal-depth one), but is retained because it is
useful metadata for the benchmark generator's stratification by circuit
depth (\cref{ch:benchmarking}).

\section{Load and Delay Computation}
\label{sec:sta-load}

\Cref{eq:cload,eq:drise,eq:dfall} are evaluated for every gate in a single
pass, independent of iteration order, since both quantities depend only on
the currently-assigned cell sizes of a gate and its fan-out --- never on an
arrival or required time. Concretely, this computation lives entirely in
\code{domain/circuit.py} and runs once, eagerly, inside
\code{Circuit.\_\_init\_\_}, via two private methods:

\begin{lstlisting}[style=pythonstyle, numbers=none]
def _compute_fanout_capacitances(self) -> dict[str, float]: ...
    # Eq. (2.1)
def _compute_gate_delays(self) -> dict[str, RiseFall]: ...
    # Eq. (2.2)-(2.3)
\end{lstlisting}

with the per-gate rise/fall delay itself computed by the module-level
helper \code{\_gate\_delay( cell: Cell, load\_capacitance: float, kappa:
float) -> RiseFall}, and total area and power by
\code{Circuit.\_compute\_area()} and \code{Circuit.\_compute\_power()}
(\cref{ch:power-area}). This is a common point of confusion worth stating
explicitly: delay is \emph{not} recomputed as a function of arrival
time during the forward pass; it is a fixed property of the current sizing
assignment, computed once by \code{Circuit.\_\_init\_\_} and only ever
\emph{read} by \code{analysis/sta.py}, never recomputed there.

\section{Forward Propagation}
\label{sec:sta-forward}

\Cref{eq:forward-sta} is evaluated for every gate in topological order by
\code{\_compute\_arrival\_times} (\cref{sec:sta-graph}), which delegates
the per-gate, per-transition maximization to
\code{\_select\_pred \newline ecessors(circuit, gate, output\_transition,
arrival\_times, gate\_delays) ->\newline tuple[float, tuple[PredecessorArc, ...]]}.
Two implementation details, both visible directly in this function's body,
are worth making explicit.

\paragraph{Net-keyed timing tables.} Arrival times are stored in a single
\code{dict[str, RiseFall]} keyed by \emph{net} name, not by gate name. A
primary input's net and a gate's output net are both, from this table's
point of view, just nets with an associated
$(\mathrm{AT}^{\mathrm{rise}}, \mathrm{AT}^{\mathrm{fall}})$
pair; primary-input nets have theirs supplied directly from
\code{circuit.config.input\_arrival(net \_name)}, while gate-output nets have
theirs computed by \cref{eq:forward-sta}. This avoids the need for a
downstream consumer of an arrival time to first determine ``is this a
primary input or a gate output'' before looking it up.

\paragraph{Exhaustive tie recording.} \code{\_select\_predecessors} builds a
list of every \code{(candidate\_arrival, (pin\_number,
input\_transition))} pair across a gate's input pins and their legal input
transitions (\code{\_valid\_input\_transitions}, implementing
\cref{sec:bg-unate}), takes \code{maximum\_arrival = max(...)}, and returns
\emph{every} candidate whose arrival exactly equals that maximum as a tuple
of \code{PredecessorArc}s --- not merely the first one encountered. This
matters directly for critical-path reconstruction (\cref{sec:sta-paths}): a
circuit with a genuinely symmetric or reconvergent structure can have two
structurally distinct paths that are exactly equally critical, and
discarding all but one of them arbitrarily would silently under-report the
set of paths actually responsible for the reported worst-case delay.

\section{Backward Propagation}
\label{sec:sta-backward}

\Cref{eq:backward-sta} is evaluated for every gate in \emph{reverse}
topological order by \code{\_compute\_required \_times}, seeded at the
primary outputs with \code{circuit.config.output\_required(net\_name)}. The
propagation of a single required time backward through one timing arc is
isolated into its own pure function,

\begin{lstlisting}[style=pythonstyle, numbers=none]
def _propagate_required(output_required: RiseFall, gate_delay: RiseFall,
                        timing_sense: TimingSense) -> RiseFall: ...
\end{lstlisting}

whose three branches implement \cref{sec:bg-unate}'s three timing-sense
cases directly: positive-unate subtracts $D_i^{s}$ from
$\mathrm{RT}_i^{s}$ transition-for-transition; negative-unate does the same
with rise and fall swapped; and, for a non-unate pin, both output
transitions are legal regardless of which input transition is being
evaluated, so the function computes
\code{earliest\_candidate = min(output\_required.rise - gate\_delay.rise,
output\_required.fall - gate\_d\newline elay.fall)} once and returns it as
\emph{both} the rise and fall required time propagated through that pin --- a
conservative choice that is required for correctness, since a non-unate
input transition's required time must respect whichever output transition
it might actually cause. \code{\_compute\_required\_times} then folds this
per-arc result into the running required time of each input net via a
plain \code{min(...)}, implementing the outer minimization of
\cref{eq:backward-sta} over every fan-out arc.

\section{Slack, WNS, and TNS}
\label{sec:sta-slack}

\Cref{eq:slack,eq:wns-tns} are evaluated by
\code{\_compute\_transition\_slacks} and
\code{\_compute\_o\newline utput\_slack\_metrics} directly from the arrival- and
required-time tables produced above. A configuration flag,
\code{config.timing\_analysis.separate\_rise\_fall}
(\code{TimingAnalysisConfig}, \cref{tab:arch-domain-api}), read inside
\code{\_output\_slack\_values}, controls whether rise and fall slack at a given
primary output are treated as two independent entries in the WNS/TNS
computation (the default, and the behavior implied by a literal reading of
\cref{eq:wns-tns}) or collapsed to a single per-output worst-case value
first; this is useful when a downstream consumer only cares
about the offending net, not which polarity of transition is responsible.

\section{Critical-Path Extraction}
\label{sec:sta-paths}

A na\"ive definition of ``the $K$ worst paths'' as ``the $K$ worst
(endpoint, transition) pairs, each with one arbitrarily chosen predecessor
chain'' would silently under-represent circuits with tied critical paths.
Path extraction is instead implemented by \code{\_find\_critical\_paths},
which iterates every primary-output net and both of its transitions and
calls the recursive helper:

\begin{lstlisting}[style=pythonstyle, numbers=none]
def _trace_paths_backward(circuit, current_net, current_transition,
        output_net, reversed_steps, reversed_transitions, endpoint_slack, found_paths, predecessors) -> None: ...
\end{lstlisting}

which \emph{branches} at every tie recorded by \code{\_select\_predecessors}
(\cref{sec:sta-forward}): it looks up \code{predecessors[(current\_net.name,
current\_transition)]}, and recurses once per stored \code{PredecessorArc}
in that tuple, prepending a \code{PathStep(gate, input\_pin)} to
\code{rever\newline sed\_steps} each time, until \code{current\_net.net\_type is
NetType.INPUT}, at which point a complete \code{CriticalPath(path, transitions,
slack)} is appended to \code{found\_paths}. All paths produced by this
search, across all endpoints, are pooled into a single list, sorted by
\code{path.slack}, and truncated to
\code{circuit.config.timing\_analysis.top\_k\_paths} by
\code{\_find\_critical\_paths} itself. Two consequences of this design are
worth stating explicitly, since they differ from a simpler ``one path per
worst endpoint'' scheme: first, the reported top-$K$ list is not guaranteed
to contain one path per distinct endpoint --- a circuit with several tied
routings to the same endpoint can occupy multiple slots in the list; second,
this is a strictly more complete answer to ``which paths are responsible
for the worst delay'' than an implementation that picks one arbitrary
routing per endpoint, at the cost of a somewhat less predictable
relationship between $K$ and the number of distinct endpoints represented.

\section{Worked Example: \texttt{c01\_inverter\_chain}}
\label{sec:sta-worked-example}

To ground the equations of this chapter in concrete numbers, this section
traces the engine's computation by hand on the smallest supplied benchmark
circuit, a three-gate inverter--buffer--inverter chain:

\begin{lstlisting}[language=bash, style=pythonstyle, numbers=none]
INPUT A
G1 INV_X1 A  N1
G2 BUF_X1 N1 N2
G3 INV_X1 N2 OUT
OUTPUT OUT
\end{lstlisting}

with primary-input arrival times $\mathrm{AT}_A^{\mathrm{rise}} =
\mathrm{AT}_A^{\mathrm{fall}} = 0$~ns, primary-output required times
$\mathrm{RT}_{\mathrm{OUT}}^{\mathrm{rise}} =
\mathrm{RT}_{\mathrm{OUT}}^{\mathrm{fall}} = 0.11$~ns, and an external load
of $4.0$~fF on \code{OUT}. \Cref{tab:c01-worked} reproduces, gate by gate,
the load, delay, arrival time, required time, and slack computed by
\cref{eq:cload,eq:drise,eq:dfall,eq:forward-sta,eq:backward-sta,eq:slack}
using the \code{INV\_X1}/\code{BUF\_X1} library entries. The resulting
$\mathrm{WNS} = 0.009$~ns and $\mathrm{TNS} = 0$~ns (no violation) were
confirmed to match the tool's reported output to six decimal places.

\begin{table}[H]
\centering
\small
\begin{tabular}{@{}lccccc@{}}
\toprule
\textbf{Gate} & \textbf{$C_{\mathrm{load}}$ (fF)} & \textbf{$D_{\mathrm{rise}}/D_{\mathrm{fall}}$ (ns)} & \textbf{$\mathrm{AT}^{\mathrm{rise}}/\mathrm{AT}^{\mathrm{fall}}$} & \textbf{$\mathrm{RT}^{\mathrm{rise}}/\mathrm{RT}^{\mathrm{fall}}$} & \textbf{Slack (ns)} \\
\midrule
G1 (\code{INV\_X1}) & 1.0 & 0.0240 / 0.0185 & 0.0240 / 0.0185 & 0.0350 / 0.0275 & 0.011 / 0.009 \\
G2 (\code{BUF\_X1}) & 1.0 & 0.0370 / 0.0315 & 0.0610 / 0.0500 & 0.0720 / 0.0590 & 0.011 / 0.009 \\
G3 (\code{INV\_X1}) & 4.0 & 0.0510 / 0.0380 & 0.1010 / 0.0990 & 0.1100 / 0.1100 & 0.009 / 0.011 \\
\bottomrule
\end{tabular}
\caption{Hand-verified static timing analysis of
\texttt{c01\_inverter\_chain}. G3 is the sole primary-output-driving gate;
its rise slack of \SI{0.009}{ns} equals the circuit-level WNS.}
\label{tab:c01-worked}
\end{table}

Two points are worth drawing attention to that a reader might otherwise
gloss over. First, because \code{INV\_X1} is negative-unate, G1's rise
\emph{output} arrival time is driven by the primary input's \emph{fall}
transition ($\mathrm{AT}_A^{\mathrm{fall}} = 0 \Rightarrow
\mathrm{AT}_{G1}^{\mathrm{rise}} = 0 + D_{\mathrm{rise},G1} = 0.024$), not
its rise transition; a reader tracing the table without accounting for
timing sense would expect the opposite pairing. Second, the worst slack in
the circuit (\SI{0.009}{ns}) alternates which transition it belongs to at
every stage of the chain ($\mathrm{fall}$ at G1, $\mathrm{fall}$ at G2,
$\mathrm{rise}$ at G3) precisely because of this same negative-unate
alternation --- a further illustration of why rise and fall must be
propagated as genuinely independent quantities rather than a single
worst-case delay per gate.
